% ===================================================================
%  نقشہ نمبر 4 — بھری جوانی تے بُڈھاپا۔
%  Traduction 2026 d'apres texto/io, controlee sur texto/fr, texto/en
%  et texto/hi. Consigne : voir texto/pnb/10-naqsha-01.tex.
%
%  LA PARENTE SE NOMME PAR LE COTE. Le pendjabi n'a pas un mot pour
%  « beau-pere » mais deux, selon qu'on parle du pere du mari ou de
%  celui de la femme ; de meme pour la belle-sœur — نݨان est la sœur
%  du mari. On suit l'ido, qui dit toujours de quel cote il parle.
% ===================================================================

%%K t04-apar-1 apar
\VUsaut{4.0mm}
\VUtitre{15.0pt}{60}{نقشہ نمبر 4}
\VUsaut{1.6mm}
\VUtitre{11.4pt}{0}{بھری جوانی تے بُڈھاپا۔}
\VUsaut{3.2mm}

%%K t04-tit-1 sub
\VUcentre{11.4pt}{0}{\textit{پہلا منظر۔}}
\VUsaut{1.6mm}
\VUtitre{11.4pt}{0}{ویاہ دی دعوت۔}
\VUsaut{2.6mm}

%%K t04-tit-2 sub
\VUpk{10.2pt}{40}{پتجھڑ۔}
\VUsaut{3.0mm}

%%K t04-c1-01-1 p
1. --- اوہ جوان ہُݨ وڈے ہو گئے نیں۔ اوہناں وچوں اک، جان، ویاہ کر رہیا اے۔ اسیں
\VUgras{ویاہ دی دعوت} وچ آں، جیہڑی لاڑے تے لاڑی دے دوہاں ٹبراں نوں اکو میز دے
چوہاں پاسیں لے آندی اے۔

%%K t04-c1-02-1 p
2. --- اسیں \VUgras{پتجھڑ} وچ آں، \VUgras{ستمبر} دے مہینے وچ۔ \VUgras{اکتوبر}
تے \VUgras{نومبر} دی ٹھری ہوا نے ہالے پنڈ توں اوہدے ہرے پتّر نئیں کھوہے۔ شام
دی ہوا ہولی تے مہکدی اے؛ ایسے لئی کھاݨا اوس \VUgras{برامدے} \textsuperscript{(8)}
دے تھلے وَرتایا جا رہیا اے جیہڑا باغ ول کھلدا اے۔

%%K t04-c1-03-1 p
3. --- دور \VUgras{ٹبی} \textsuperscript{(3)} اُتے جان دے ماں پیو دا گھر اے، اک
سجیلا \VUgras{پنڈ دا مکان} \textsuperscript{(1)}، ہریاول وچ لُکیا ہویا؛ چنگے
انگور دے باغ، جیہناں توں ودھیا شراب بݨدی اے، ڈھلوان نوں ڈھکے ہوئے نیں۔ انگور
اُگاؤݨ والا اوہناں دی سنبھال کردا اے۔

%%K t04-c1-04-1 p
4. --- ویلاں دے اُتوں \VUgras{تِتراں} \textsuperscript{(2)} دا اک جھُنڈ اُڈدا اے،
جیہڑا \VUgras{شکار دے رکھوالے} \textsuperscript{(5)} دے نیڑے آوݨ نال چونک گیا اے۔
اوہ، بغل وچ \VUgras{بندوق} دبائی تے \VUgras{مونڈھے اُتے شکار دا تھیلا} لٹکائی،
\VUgras{جھاڑیاں} \textsuperscript{(4)} کُٹ رہیا اے، اپݨے \VUgras{کُتے}
\textsuperscript{(6)} نال۔ اوہ جاگیر اُتے نظر رکھدا اے تے چور شکاریاں نوں شکار
تباہ نئیں کرن دیندا۔

%%K t04-c1-05-1 p
5. --- \VUgras{برامدے} دے باہر اک \VUgras{سِدھائی ہوئی ویل} چڑھدی اے، جیدے توں
بھارے \VUgras{انگوراں دے گُچھے} \textsuperscript{(7)} لٹکدے نیں۔

%%K t04-c1-06-1 p
6. --- \VUgras{میز} \textsuperscript{(16)} نوں سجاؤݨ والے پتجھڑ دے سوہݨے پھُل
\VUgras{گُل داؤدی} \textsuperscript{(9)} نیں؛ لاڑی دے \VUgras{پیو}
\textsuperscript{(10)} نے اپݨے کوٹ دے کاج وچ اک لایا ہویا اے۔

%%K t04-c1-07-1 p
7. --- کھاݨا مُکݨ اُتے اے۔ \VUgras{نوکر} \textsuperscript{(11)} مِٹھے دے برتن
لیاؤݨ لگ پیا اے تے تھوڑی دیر مگروں اوہ سب کجھ چُک لے جاوے گا جیدی ہُݨ لوڑ نئیں
— \VUgras{چمچے} \textsuperscript{(14)}، \VUgras{کانٹے} \textsuperscript{(12)}،
\VUgras{چھُریاں} \textsuperscript{(13)}، \VUgras{وَرتاؤݨ دے برتن}
\textsuperscript{(15)}۔

%%K t04-tit-3 sub
\VUpk{10.2pt}{40}{ٹبر۔}
\VUsaut{3.0mm}

%%K t04-c1-08-1 p
8. --- سارے خوش وکھائی دیندے نیں، تے لاڑے دی \VUgras{ماں} \textsuperscript{(17)}
جیس مُسکان نال اپݨے نِکیاں نوں ویکھ رہی اے، اوہ دسدی اے پئی اوہ کِنّی سُکھی اے۔

%%K t04-c1-09-1 p
9. --- ایہہ ویاہ علاقے دے دو کھاندے پیندے ٹبراں نوں جوڑدا اے۔ جان،
\VUgras{گھر والا} \textsuperscript{(18)}، اک وڈے زمیندار دا \VUgras{پُتر} اے، تے
اوہ اک کامیاب کارخانے دار دا \VUgras{جوائی} بݨ رہیا اے۔

%%K t04-c1-10-1 p
10. --- \VUgras{جوڑی} جوان اے، فیر وی اوہناں دی \VUgras{کُڑمائی} بہت پہلاں ہو
چُکی سی۔ \VUgras{نویں لاڑی} \textsuperscript{(19)}، مارتھا، سنترے دے پھُلاں نال
سجی اپݨی \VUgras{چٹی پوشاک} وچ موہ لیݨ والی لگدی اے۔

%%K t04-c1-11-1 p
11. --- اوہدے \VUgras{سوہرا} \textsuperscript{(20)} اِنّی چنگی \VUgras{نونہہ}
لبھ کے خوش نیں، تے جان دی \VUgras{سس} \textsuperscript{(21)} اپݨی \VUgras{دھی}
نوں اِنّی سوہݨی ویکھ کے مُسکرا رہی اے۔

%%K t04-c1-12-1 p
12. --- میز دے سرے اُتے دوجے \VUgras{پراہُݨے} \textsuperscript{(22)} بیٹھے نیں:
\VUgras{رشتے دار، یار، چچیرے بھرا}۔ اک \VUgras{مکھ ویٹر} \textsuperscript{(23)}،
بانہہ اُتے رومال پائی، پھرتی نال اوہناں دی خدمت کر رہیا اے۔

%%K t04-tit-4 sub
\VUpk{10.2pt}{40}{یار۔}
\VUsaut{3.0mm}

%%K t04-c1-13-1 p
13. --- میز دے سجے پاسے ایہہ رہی اک \VUgras{مُٹیار} \textsuperscript{(24)}، لاڑی
دی \VUgras{سہیلی}، فیر لاڑی دا \VUgras{بھرا}، جیہڑا اوہدا \VUgras{سربالا}
\textsuperscript{(25)} اے۔ اوہ اپݨی \VUgras{سربالی} \textsuperscript{(26)} نال گلاں
کر رہیا اے، جیہڑی لاڑے دی بھین اے تے ایس ویاہ نال مارتھا دی \VUgras{نݨان} بݨ
رہی اے۔ اوہ موہ لیݨ والی مُٹیار اے۔ اوہدے کول سوہݨے \VUgras{گہݨے} نیں، اک
\VUgras{موتیاں دا ہار} تے سونے دا اک \VUgras{کنگݨ}۔

%%K t04-c1-14-1 p
14. --- اوہناں کول کھڑے ہو کے گلاس چُکدے صاحب ٹبر دے اک \VUgras{یار}
\textsuperscript{(27)} نیں۔ اوہ \VUgras{لاڑے دے گواہاں} وچوں اک نیں۔ اوہ
\VUgras{وَدھائی} دے رہے نیں، نویں ویاہیاں دی \VUgras{صحت} دی کامنا کر رہے نیں تے
اوہناں نوں لمے سُکھ دی دعا دے رہے نیں۔ اوہ پورے ٹھاٹھ وچ نیں: لما کوٹ تے
\VUgras{چمڑے دے چمکدے جُتے} \textsuperscript{(28)}۔ اوہ \VUgras{دادی}
\textsuperscript{(29)} نوں نال لیائے نیں، جیہڑی \VUgras{رنڈی} اے۔

%%K t04-c1-15-1 p
15. --- \VUgras{لمے کوٹ} \textsuperscript{(31)} والا اوہ جوان، جیدی پتلی کالی مُچھ
اے، جان دا \VUgras{بھݨوئیا} \textsuperscript{(30)} اے۔ اوہ مشہور انجینیر اے۔ اوہ
اک بہت سجیلی \VUgras{مُٹیار} \textsuperscript{(33)} کول بیٹھا اے، جیہڑی مارتھا دی
\VUgras{وڈی بھین} اے تے اوس پیاری نِکی کُڑی دی ماں، جیہڑی اپݨے چچیرے بھرا دا ہتھ
پھڑی پھُل دیݨ تے \VUgras{« شام سلام »} \VUgras{آکھݨ} آ رہی اے۔ ایس بیبی دے
\VUgras{کن پھُل} بہت سوہݨے نیں تے اوہدے \VUgras{والاں دے سنگار} وچ سلیقہ اے۔

%%K t04-c1-16-1 p
16. --- اوہ نِکا چچیرا بھرا، جیہڑا اپݨی \VUgras{کُڑتی} \textsuperscript{(36)} تے
اپݨی ڈھِلی \VUgras{نِکر} \textsuperscript{(37)} وچ اِنّا انوکھا لگدا اے، بڑا ترس
جوگ اے۔ اوہ \VUgras{یتیم} \textsuperscript{(35)} اے۔ بہت نِکی عمر وچ اوہنے اپݨا
پیو تے اپݨی ماں گنوا لئی۔ اوہدے چاچا اوہدے \VUgras{سرپرست} \textsuperscript{(38)}
نیں، تے اوس وچارے بال نوں پالݨ لئی ای چھڑے رہ گئے۔ اوہ بھلے بندے نیں۔ اوہناں دے
سر اُتے وال گھٹ نیں تے اوہناں نوں دور دا بہت گھٹ دِسدا اے؛ اوہ \VUgras{نک عینک}
لاندے نیں۔

%%K t04-tit-5 sub
\VUsaut{2.4mm}
\VUpk{10.2pt}{40}{مِٹھا۔}
\VUsaut{3.0mm}

%%K t04-c1-17-1 p
17. --- پراہُݨیاں دے پچھے، \VUgras{میز پوش} وِچھی اک میز اُتے، \VUgras{مِٹھا}
\textsuperscript{(39)} سجیا ہویا اے۔ اک وڈے کٹورے وچ ایہہ رہے پھل:
\VUgras{آڑو} \textsuperscript{(40)}، \VUgras{ناشپاتی} \textsuperscript{(41)}،
\VUgras{سیب} \textsuperscript{(42)}، \VUgras{خوبانی} \textsuperscript{(43)} تے
\VUgras{انگور} \textsuperscript{(44)}۔ اوہناں کول \VUgras{انناس}
\textsuperscript{(45)}، اک سوہݨا \VUgras{کیک} \textsuperscript{(46)}،
\VUgras{شراب دیاں بوتلاں} \textsuperscript{(48)} تے \VUgras{شیمپین}
\textsuperscript{(49)}، تے صاف \VUgras{پلیٹاں} \textsuperscript{(47)} دے ڈھیر وی
نیں۔

%%K t04-c1-18-1 p
18. --- \VUgras{شراب وَرتاؤݨ والا} \textsuperscript{(50)} پرانی شراب دی اک
\VUgras{بوتل} \textsuperscript{(52)} کھول رہیا اے، \VUgras{پیچ}
\textsuperscript{(51)} نال۔ \VUgras{ڈاٹ} \textsuperscript{(53)} بڑی اوکھ نال نکلدی
لگدی اے۔ ایس وَرتاؤݨ والے دی بانہہ اُتے اک \VUgras{رومال} \textsuperscript{(54)}
اے۔

%%K t04-c1-19-1 p
19. --- کھاݨے مگروں صاحب لوک سگرٹ والے کمرے وچ جاݨ گے، تے بیبیاں ناچ لئی تیار
ہوݨ ٹُر پیݨ گیاں۔

%%K t04-tit-6 sub
\VUsaut{5.0mm}
\VUcentre{11.4pt}{0}{\textit{دوجا منظر۔}}
\VUsaut{1.6mm}
\VUpk{11.4pt}{40}{دادے دا جنم دن۔}
\VUsaut{2.6mm}

%%K t04-tit-7 sub
\VUpk{10.2pt}{40}{سوغات۔}
\VUsaut{3.0mm}

%%K t04-c2-01-1 p
1. --- دس ورھے لنگھ گئے۔ جان دے ماں پیو بُڈھے ہو چلے نیں، تے اوہ اپݨے تِناں
پوتریاں نوں بہت پیار کردے نیں — دو \VUgras{پوترے} \textsuperscript{(2)} تے اک
\VUgras{پوتری} \textsuperscript{(3)}۔ اج \VUgras{دادے دا جنم دن} اے، ایسے لئی
نِکے اوہناں نوں وَدھائی دیݨ آئے نیں۔

%%K t04-c2-02-1 p
2. --- نِکا پیٹر ہالے بہت نِکا اے؛ اوہ نِرا تن ورھیاں دا اے۔ اوہ \VUgras{بلی}
\textsuperscript{(1)} نوں اپݨے ول آندی ویکھدا اے۔ اوہ اوہدی سوہݨی ریشمی
\VUgras{جھلّ} تے لمی \VUgras{پُوچھ} اُتے ہتھ پھیرنا چاہندا اے؛ اوہنوں ایہہ نئیں
سُجھدا پئی اوہناں سوہݨے \VUgras{پنجیاں} دے تھلے تِکھے نوہہ نیں، جیہڑے اوہنوں
لِتاڑ وی سکدے نیں۔

%%K t04-c2-03-1 p
3. --- اوہدی بھین گیبریئل، جیہڑی اوہدا ہتھ پھڑی ہوئی اے، کِتے ودھ سنجیدہ اے، تے
مارسل، اپݨے وڈے \VUgras{اوورکوٹ} \textsuperscript{(4)} تے چمڑے دی \VUgras{پیٹی}
\textsuperscript{(5)} وچ، پورا نِکا بندہ لگدا اے، بھانویں نِکی نِکر اوہدے
\VUgras{جُرابا} \textsuperscript{(6)} کھلے چھڈ دیندی اے۔

%%K t04-c2-04-1 p
4. --- اوہناں دے پیو نے اپݨا موٹا سیال دا \VUgras{اوورکوٹ} \textsuperscript{(7)}
پایا ہویا اے، \VUgras{جھلّ والے کالر} \textsuperscript{(8)} سمیت، تے اپݨی
\VUgras{جیب} \textsuperscript{(9)} وچ اک گلوبند رکھیا ہویا اے۔ اوہدی گھر والی نے
\VUgras{پراں} \textsuperscript{(10)} نال سجی نمدے دی ہیٹ، اپݨی \VUgras{جیکٹ}
\textsuperscript{(11)}، اپݨا جھلّ والا \VUgras{گلوبند} \textsuperscript{(12)} تے
اپݨا \VUgras{مف} \textsuperscript{(13)} پایا ہویا اے۔

%%K t04-c2-05-1 p
5. --- بہت پالا اے، کیوں جے سیال آ گیا اے۔ \VUgras{برف} \textsuperscript{(19)}
سارا دن پیندی رہی اے تے گھراں دیاں چھتاں بالکل چٹیاں نیں۔ \VUgras{رات} نال پالا
ہور وی تِکھا ہو جاندا اے۔ اسمان وچ \VUgras{چن} \textsuperscript{(17)} تے
\VUgras{تارے} \textsuperscript{(18)} چمکدے نیں تے \VUgras{ہنیر} نوں چیردے نیں۔

%%K t04-tit-8 sub
\VUsaut{4.2mm}
\VUpk{10.2pt}{40}{بیماری۔}
\VUsaut{3.0mm}

%%K t04-c2-06-1 p
6. --- اوہ وچارا \VUgras{انھا} \textsuperscript{(20)}، جیہڑا \VUgras{دواخانے}
\textsuperscript{(16)} دے سامݨے چونک لنگھ رہیا اے تے جیہنوں اوہدا \VUgras{پوڈل}
\textsuperscript{(21)} راہ وکھا رہیا اے، بڑا نِگھرا اے۔ \VUgras{نوکراݨی}
\textsuperscript{(14)} جسٹین رسوئی توں بہت گرم \VUgras{کاڑھے} دا اک
\VUgras{کٹورا} \textsuperscript{(15)} لیا رہی اے، خوب مِٹھا کیتا ہویا۔

%%K t04-c2-07-1 p
7. --- \VUgras{دادی} \textsuperscript{(23)}، جیہڑی میز اُتوں اپݨی
\VUgras{عینک} \textsuperscript{(26)} چُکݨا چاہندی اے تاں جو اوہ \VUgras{پرچی}
\textsuperscript{(27)} پڑھ سکے جیہڑی \VUgras{حکیم} \textsuperscript{(25)} نے ہُݨے
لکھی اے، اپݨی \VUgras{سوٹی} \textsuperscript{(24)} دے سہارے اپݨی کرسی توں اُٹھ
رہی اے۔

%%K t04-c2-08-1 p
8. --- اوہناں نوں اکثر نِکیاں نِکیاں \VUgras{تکلیفاں} رہندیاں نیں،
\VUgras{چکر}، \VUgras{زکام} تے \VUgras{دند دی پیڑ}؛ اوہناں دیاں لتاں کمزور نیں
تے اوہناں دیاں اکھاں ہُݨ چنگیاں نئیں رہیاں۔ فیر وی اوہ اپݨے بُڈھے
\VUgras{حکیم} نوں چھیڑ کے خوش ہوندیاں نیں، جیہڑے اوہناں نوں ہمیشہ کوئی
\VUgras{گھُٹ} \textsuperscript{(28)} لکھ دیݨا چاہندے نیں۔ اج اوہ اپݨا جوان ٹبر
اپݨے چوہاں پاسیں لبھ کے بہت سُکھی نیں۔

%%K t04-c2-09-1 p
9. --- چنگی طرحاں بند کمرے وچ نیم اے۔ سنگ مرمر دی \VUgras{انگیٹھی}
\textsuperscript{(29)} وچ \VUgras{سوہݨی اگ} \textsuperscript{(30)} بل رہی اے، تے
ہُݨے ہُݨے بالے \VUgras{لیمپ} \textsuperscript{(31)} دی کُونی \VUgras{لو} وچ دل
خوش ہو جاندا اے۔

%%K t04-tit-9 sub
\VUsaut{4.2mm}
\VUpk{10.2pt}{40}{دادا۔}
\VUsaut{3.0mm}

%%K t04-c2-10-1 p
10. --- \VUgras{دادا} \textsuperscript{(33)} تیکر بھُل گئے نیں پئی اوہ بیمار رہے
نیں، پئی اوہناں دا \VUgras{جوڑاں دا درد} اوہناں نوں اپݨی \VUgras{آرام کرسی}
\textsuperscript{(34)} وچ بنھی رکھدا اے، تے پئی کل ای اوہناں نوں
\VUgras{تاپ تے غشی} آئی سی۔ اوہناں نوں ہُݨ یاد نئیں پئی پچھلے سیال اوہناں نوں
\VUgras{نمونیا} ہویا سی، تے اک سخت پیڑ دیݨ والی \VUgras{کھنگ} نے اوہناں نوں بہت
ستایا سی۔

%%K t04-c2-10-2 p
تے اوہ اوکھ نال ای راضی ہوئے۔ ہُݨ کجھ چنگے نیں۔ پر اوہناں دی لت، جیہنوں اوہ اک
\VUgras{گدی} \textsuperscript{(38)} اُتے ٹِکائی رکھدے نیں، تے اوہ گدی اک نِکی
\VUgras{چوکی} \textsuperscript{(39)} اُتے رکھی اے، ہُݨ وی دُکھدی اے۔

%%K t04-c2-11-1 p
11. --- جدوں اوہ اپݨے نیمے \VUgras{ڈریسنگ گاؤن} \textsuperscript{(35)} وچ، اپݨے
وڈے سِپی دے \VUgras{بٹناں} \textsuperscript{(36)} سمیت، چنگی طرحاں لپیٹے ہوندے
نیں، تے جدوں اوہناں کول اخبار پڑھݨ لئی اوہ نِکی \VUgras{یتیم کُڑی}
\textsuperscript{(40)} ہوندی اے، اپݨی چٹی \VUgras{ٹوپی}، اپݨی نیلی
\VUgras{پوشاک} \textsuperscript{(42)} تے \VUgras{چادر} \textsuperscript{(41)} وچ،
تدوں اوہ کوئی شکایت نئیں کردے۔

%%K t04-c2-12-1 p
12. --- ہر ورھے، جدوں \VUgras{جنتری} \textsuperscript{(43)} پہلی \VUgras{دسمبر}
دی \VUgras{تریخ} مُڑ لے آندی اے، اوہ اپݨے \VUgras{پوتریاں} دی بے صبری نال اُڈیک
کردے نیں، کیوں جے اوہ جاݨدے نیں پئی اوہ اوس اہم \VUgras{تریخ} نوں نئیں بھُلݨ
گے۔

%%K t04-c2-13-1 p
13. --- اوہ جذباتی ہو کے اوہ دور دے دن یاد کردے نیں جدوں اوہ آپ اوس شاندار شاہی
\VUgras{مارشل} نوں وَدھائی دیݨ جایا کردے سن، جیہناں دی \VUgras{تصویر}
\textsuperscript{(44)} کندھ اُتے ٹنگی اے، تے جیہڑے نِکیاں دے \VUgras{پڑ دادا} نیں۔
\VUsaut{6.0mm}
\VUfilet{22mm}
